% Generated draft: requires mathematical and source review.
\documentclass[11pt]{article}
\usepackage[margin=25mm]{geometry}
\usepackage{fontspec}
\setmainfont{Latin Modern Roman}
\usepackage{amsmath,amssymb,mathtools,unicode-math}
\setmathfont{Latin Modern Math}
\usepackage{enumitem,needspace,xurl,hyperref}
\setlength{\parindent}{0pt}
\setlength{\parskip}{5pt}
\setlength{\emergencystretch}{3em}
\newcommand{\N}{\mathbb{N}}
\newcommand{\Z}{\mathbb{Z}}
\newcommand{\Q}{\mathbb{Q}}
\newcommand{\R}{\mathbb{R}}
\newcommand{\C}{\mathbb{C}}
\newcommand{\F}{\mathbb{F}}
\newcommand{\A}{\mathbb{A}}
\newcommand{\PP}{\mathbb{P}}
\newcommand{\E}{\mathbb{E}}
\newcommand{\eps}{\varepsilon}
\newcommand{\dd}{\,\mathrm{d}}
\begin{document}
\section*{UnsolvedMath ID: 30004297}
\textbf{Source problem code:} catalogue entry
\subsection{Definitions and mathematical statement}
For a real number \(B\geq 1\), define the set of integer triples
\[
S(B)=\{(x,y,z)\in\Z^3: \max\{|x|,|y|,|z|\}\leq B,\quad zy^2+yz^2=x^3-xz^2\}.
\]
Using the usual interpretation of a sum of two squares, put
\[
r_0(n)=\begin{cases}
1,&\text{if there exist }a,b\in\Z\text{ such that }n=a^2+b^2,\\
0,&\text{otherwise},
\end{cases}
\qquad(n\in\Z).
\]
Thus this convention includes \(r_0(0)=1\). Interpreting the supplied sum as a sum over all the stated integer triples, define
\[
N_{\mathrm{lit}}(B)=\sum_{(x,y,z)\in S(B)}r_0(z).
\]
The literal question supplied in the registry is whether
\[
\exists C\in\R_{\geq0}\;\exists B_0\geq1\;\forall B\geq B_0,
\qquad N_{\mathrm{lit}}(B)\leq C.
\]
\textbf{Unresolved formulation qualifications.} The record does not require \(z\neq0\), primitive triples, or identification of triples representing the same projective point. The equation is homogeneous, and the literal summation includes triples with \(z=0\) and common multiples of solutions. These features must be reconciled with the intended counting problem before adopting this as a precise elliptic-curve boundedness conjecture. No such restrictions are imposed in this draft.

The supplied title and dated curation describe a rank-one elliptic curve and report an \(o(\sqrt{\log B})\) estimate for a weighted integral-point count. The record does not provide the hypotheses or counting conventions needed to identify that count with \(N_{\mathrm{lit}}(B)\). The supplied open-status label and literature assessment are recorded as source claims, not independently verified mathematical conclusions.
\subsection{Short English statement}
The supplied question asks whether the number of bounded integer solutions of the displayed equation, weighted by whether \(z\) is a sum of two squares, stays bounded as the height bound grows. The literal triple count is defined above, but missing restrictions concerning zero coordinates, common factors, and projective representatives prevent a reliable identification with the intended rank-one elliptic-curve problem.
\subsection{Sources}
\begin{itemize}
\item Supplied local UnsolvedMath registry record, permanent ID \(30004297\), source problem code \texttt{catalogue entry}, in the Oberwolfach Reports Open Problems collection. The record supplies no named proposer and lists the proposed year as 2019.
\item Supplied report citation: \emph{Analytic Number Theory (2020)}, Oberwolfach Reports, pp.\ 3141--3205, DOI \texttt{10.4171/owr/2019/50}. Supplied source URL and literature-source link, labelled ``original OWR report'': \url{https://doi.org/10.4171/owr/2019/50}. The background describes the reports as documenting mathematical workshops and their open-problem discussions. This report was not independently consulted for this draft.
\item Literature assessment dated exactly \texttt{2026-08-22}: supplied dated curation, not independently verified. It labels the problem open, reports an \(o(\sqrt{\log B})\) estimate for a weighted integral-point count on a rank-one elliptic curve, and states that no general boundedness theorem for this count is established. The supplied background explicitly requests independent verification of time-sensitive status and cited claims.
\item The registry labels the statement \texttt{corrected\_verified}, while its verification note specifies only repairs to formula delimiters, notation, grammar, and typography, with mathematical content and scope preserved. This metadata does not resolve the mathematical formulation issues above.
\end{itemize}
\end{document}
